\subsection{Частые ошибки}

Перед запуском симулятор проверяет программу и объясняет найденное в уведомлениях (колокольчик в заголовке редактора), а во время полёта останавливает её, если команды противоречат друг другу. Вот что он чаще всего находит:

\begin{itemize}
    \item \textbf{Команды подряд без ожидания.} \verb|ap.push(Ev.MCE_PREFLIGHT)| и сразу \verb|ap.push(Ev.MCE_TAKEOFF)|. Взлёт --- после \verb|ENGINES_STARTED| или через таймер, точка --- после \verb|TAKEOFF_COMPLETE|, посадка --- после \verb|POINT_REACHED|.
    \item \textbf{Посадка раньше точки.} Таймер короче перелёта --- дрон сядет, не долетев. Садиться по событию.
    \item \textbf{Нет \texttt{callback}.} На борту такая программа останавливается после первой команды; симулятор предупреждает.
    \item \textbf{Все точки по одному событию.} Без счётчика шагов вторая точка затирает первую. Шаблон --- \verb|nextPoint| с \verb|curr|.
    \item \textbf{Точка «в земле».} \key{Ctrl}+щелчок показывает высоту поверхности; $z$ цели должен быть выше.
    \item \textbf{Градусы вместо радиан.} \verb|ap.updateYaw(90)| --- это 90 радиан, то есть 14 оборотов. Нужно \verb|math.rad(90)|.
    \item \textbf{Таймер не остановлен.} Таймер окружности продолжает слать точки после команды посадки --- дрон не сядет. \verb|t:stop()| перед сменой этапа.
    \item \textbf{Функция вызвана вместо передачи.} \verb|Timer.callLater(1, nextPoint())| выполняет \verb|nextPoint| сразу; нужно \verb|Timer.callLater(1, nextPoint)|.
\end{itemize}

\begin{hintbox}[Где смотреть]
Терминал (вкладка \ui{Вывод}) показывает \verb|print| и сообщения автопилота о каждом переходе: предстарт, взлёт, точка достигнута. Панель \ui{Полёт} показывает текущее состояние автопилота и координаты. Если дрон «завис» --- сначала туда.
\end{hintbox}
